\chapter[Band 44: Kommutative Ringe mit Eins]{Band 44 auf einen Blick}
\noindent\textbf{Kommutative Ringe mit Eins}\par\medskip
Die Ringe mit Eins aus Band~43 werden durch Kommutativität der
Multiplikation spezialisiert.

\begin{BandUebersicht}
\textbf{Kommutativer Ring}

\UebFormel{\begin{aligned}&\CommutativeRingAlg R+\cdot01\ \leftrightarrow\\&\RingAlg R+\cdot01\land\\&\forall a,b\in R\,(a\cdot b=b\cdot a).\end{aligned}}
\UebQuellen{\FormulaRefAuto{CommutativeRingDef}[def]}
&
\textbf{Grundstruktur}
Alle allgemeinen Ringsätze aus Band~43 bleiben anwendbar.
\UebFormel{\CommutativeRingAlg R+\cdot01\vdash\RingAlg R+\cdot01.}
\UebQuellen{\FormulaRefAuto{CommutativeRingBaseRingAxiom}[ax]} \\
\addlinespace[.8ex]

\textbf{Morphismen}
Quelle und Ziel sind kommutative Ringe; \(f\) ist ein Ringhomomorphismus.
\UebFormel{\CommutativeRingHom f{R,+,\cdot,0,1}{S,\oplus,\odot,0_S,1_S}.}
\UebQuellen{\FormulaRefAuto{CommutativeRingHomDef}[def]}
&
\textbf{Isomorphismen}
Zur Homomorphie kommt die Bijektivität hinzu.
\UebFormel{\CommutativeRingIso f{R,+,\cdot,0,1}{S,\oplus,\odot,0_S,1_S}.}
\UebQuellen{\FormulaRefAuto{CommutativeRingIsoDef}[def];
\FormulaRefAuto{CommutativeRingIsoBijectiveAxiom}[ax]} \\
\addlinespace[.8ex]

\textbf{Multiplikation ganzer Zahlen}

\UebFormel{z,w\in\mathbb Z\vdash z\cdot w=w\cdot z.}
\UebQuellen{\FormulaRefAuto{IntegerMulCommutativeAxiom}[ax]}
&
\textbf{Grundbeispiel}

\UebFormel{\CommutativeRingAlg{\mathbb Z}{+}{\cdot}{\IntZero}{\IntOne}.}
\UebQuellen{\FormulaRefAuto{IntegerCommutativeRing}[thm]} \\
\addlinespace[.8ex]
\end{BandUebersicht}
